\begin{abstract}
Identifying the origin of an epidemic outbreak---the so-called patient-zero or
source detection problem---is a fundamental challenge in network epidemiology.
Given the observed infection state of a contact network at some time $T$, the
task is to rank network nodes by their likelihood of having initiated the
spreading process. While graph neural network (GNN) methods have recently
demonstrated competitive performance on this task, the majority of existing
approaches operate on static projections of the contact graph, thereby
discarding the temporal ordering of contacts that governs causal transmission
paths.

This thesis addresses that gap by implementing and systematically comparing
four GNN architectures of increasing temporal expressiveness for source
detection on temporal contact networks. The evaluated models are: (i) a static
GNN baseline that processes the time-aggregated graph; (ii) the Backtracking
Network (BN) of \cite{Ru2023}, which encodes temporal activation patterns as
edge textures; (iii) the De Bruijn Graph Neural Network (DBGNN) of
\cite{Qarkaxhija2022}, which constructs higher-order De Bruijn graphs to
represent causal walks; and (iv) a DAG-GNN applied to Temporal Event Graphs,
combining the event-graph representation with directed acyclic graph
convolutions \cite{Sterchi2025}. All models are trained in a supervised
fashion using ground-truth sources generated by Monte Carlo SIR simulations on
synthetic and real-world contact networks via an adapted version of Holme's
fast temporal SIR engine.

The thesis provides what is, to our knowledge, the first systematic
multi-architecture comparison of temporal GNN methods for epidemic source
detection. We examine three research hypotheses concerning causal path
integrity, temporal granularity convergence, and signal-to-noise resilience in
late-stage outbreaks, thereby contributing both empirical benchmarks and
architectural insights to the emerging field of learning-based network
epidemiology.
\end{abstract}
